El diseño de \cref{CAP:DESIGN} fija la estructura; esta sección recoge las partes de la
implementación en las que el detalle importa, porque un error en ellas pierde trabajo del
profesorado o produce resultados incorrectos sin que la interfaz lo muestre. Los
fragmentos de código están abreviados: se han eliminado partes no relevantes para el
argumento, pero no se ha cambiado su lógica. La renovación de la sesión de vuelo único y la
selección de la ruta más específica en la simulación ya se describieron en
\cref{SEC:DES-DATA} y no se repiten.

\subsection{Calificación con concurrencia optimista}

Las operaciones de calificación del \dfn{backend} (\texttt{PUT
/instances/\{id\}/problems/\{pid\}/grade/} y su variante por \dfn{rubrica}) aplican un
control de concurrencia optimista: además de la nueva puntuación, el cliente envía en
\texttt{old\_score} la que cree vigente, y el servidor responde \texttt{409} con el código
\texttt{VERSION\_CONFLICT} si no coincide con la almacenada. El mecanismo evita que dos
correctores se sobrescriban sin saberlo, pero exige que el cliente conozca la nota anterior,
y el contrato no ofrece al profesorado ninguna operación que devuelva las notas por problema
(hueco gap-1, \cref{SEC:AN-GAPS}).

La implementación resuelve la parte que depende del \dfn{frontend}: cada navegador recuerda
las notas que él mismo ha guardado (\texttt{useKnownGrades}) y las usa como
\texttt{old\_score}. El \cref{COD:IMP-GRADE} muestra el núcleo del componente que califica un
problema. La nota recordada solo se actualiza con la respuesta del servidor, no con el
valor tecleado, de modo que el siguiente guardado parte siempre de lo que el
\dfn{backend} ha confirmado.

\begin{lstlisting}[caption={Guardado de una nota con \texttt{old\_score}
  (\texttt{features/grading/ProblemGrader.jsx}).},label={COD:IMP-GRADE}]
const current = known.grades[problem.id]; // last grade this browser saved
const save = async (fn) => {
  setState('saving');
  try {
    const res = await fn(current?.score ?? null);   // old_score
    known.remember(problem.id, { score: res.score });
    setState('saved');
  } catch (e) {
    setState('error');
    setErrCode(e.code);                  // e.g. VERSION_CONFLICT
  }
};
save((old) => gradeManual(instanceId, problem.id, n, old));
\end{lstlisting}

La solución tiene un límite conocido: si un problema se calificó desde otro navegador, la
interfaz no puede conocer su nota y el primer guardado recibe el conflicto. En ese caso se
muestra el mensaje traducido del código de error y se avisa de que la nota total procede de
otra sesión; la resolución completa requiere una operación de lectura en el
\dfn{backend}, que se ha documentado como petición al equipo del \dfn{backend}.

\subsection{Autoguardado de las anotaciones}

La capa de tinta guarda los trazos sin que el corrector tenga que pulsar ningún botón: cada
cambio programa el guardado de la página tras 1,2 segundos sin actividad, y un cambio
posterior reinicia el temporizador. La primera versión cancelaba los temporizadores
pendientes al salir del examen, de modo que se perdían los últimos trazos si el corrector
pasaba al examen siguiente antes de ese intervalo. El fallo se detectó al contrastar la
redacción del \cref{CAP:DESIGN} con el código y se corrigió en una iteración posterior.

Al escribir la prueba de regresión apareció un segundo defecto, más grave: los
identificadores de las anotaciones ya creadas se guardaban en un objeto mutable compartido
por todas las instancias. Si la creación de una anotación seguía en curso cuando el
corrector cambiaba de examen, la respuesta escribía su identificador en el objeto ya
reiniciado para el examen siguiente, y los trazos posteriores de ese examen se acababan
guardando sobre la anotación del anterior. La corrección (\cref{COD:IMP-AUTOSAVE}) tiene dos
partes: al abandonar la instancia se \emph{ejecutan} los guardados pendientes en lugar de
cancelarlos, y cada guardado captura de forma síncrona el mapa de identificadores de su
propia instancia antes de esperar a la red.

\begin{lstlisting}[caption={Guardados pendientes al salir de un examen
  (\texttt{features/grading/useAnnotations.js}).},label={COD:IMP-AUTOSAVE}]
useEffect(() => {
  ids.current = { stylus: {}, text: null };   // fresh map per instance
  /* ... load the instance's annotations ... */
  return () => {
    const pending = Object.values(timers.current);
    timers.current = {};
    pending.forEach(({ id, fn }) => { clearTimeout(id); fn(); });  // flush, not cancel
  };
}, [instanceId]);

const persistPage = async (page) => {
  const own = ids.current;      // captured before awaiting the network
  const id = own.stylus[page];
  if (id) await updateAnnotation(id, { payload });
  else own.stylus[page] = (await createAnnotation(instanceId, body)).id;
};
\end{lstlisting}

Ninguno de los dos defectos producía un error visible: la interfaz mostraba el estado
\emph{guardado} y el problema solo se habría descubierto al reabrir el examen. Es un ejemplo
del tipo de fallo que las pruebas de flujo detectan y la revisión superficial no
(\cref{SEC:TEST-FLOWS}).

\subsection{Zonas de reconocimiento y sistemas de coordenadas}

El editor de zonas permite marcar sobre cada página del modelo de examen las regiones que el
\dfn{backend} leerá en los escaneos: el código QR, campos de texto o numéricos y casillas.
Las anotaciones del \dfn{frontend} usan fracciones del tamaño de la página con el origen en
la esquina superior izquierda, como el resto de la interfaz. El \dfn{backend}, en cambio,
almacena las zonas en puntos \acs{pdf} con el origen en la esquina inferior izquierda,
relativas a las dimensiones del perfil de página. Además, la respuesta real de la operación
de zonas es una lista donde el documento OpenAPI declara un objeto (gap-8).

La conversión se aisló en un módulo puro con sus dos sentidos (\cref{COD:IMP-ZONES}), lo que
permite probarla sin interfaz. Los valores se redondean a centésimas de punto para que una
zona guardada y vuelta a leer conserve exactamente su geometría.

\begin{lstlisting}[caption={Conversión de la caja del editor a la geometría de la
  \acs{api} (\texttt{features/exams/setup/zoneGeometry.js}).},label={COD:IMP-ZONES}]
export function toPdfBox({ x, y, width, height }, { page_width: W, page_height: H }) {
  return {
    x: r2(x * W),
    y: r2((1 - y - height) * H),   // top-left origin -> bottom-left origin
    width: r2(width * W),
    height: r2(height * H),
  };
}
\end{lstlisting}

\subsection{Reconocimiento de escritura sobre anotaciones}

El \dfn{backend} ofrece dos operaciones de \acs{ocr} sobre anotaciones de tipo trazo:
\texttt{ocr/}, que devuelve el texto reconocido, y \texttt{ocr-grade/}, que interpreta una
cifra como nota de un problema y la aplica automáticamente si la confianza supera un umbral
configurado. En la interfaz se materializan como \emph{Nota a mano}, en cada problema de
puntuación manual, y \emph{Escribir a mano}, en el comentario general.

La decisión relevante fue no reutilizar la tinta de la página. Al leer el renderizador del
\dfn{backend} se comprobó que solo dibuja listas de puntos con coordenadas de 0 a 100,
mientras que la tinta del corrector se guarda en un formato propio versionado con
coordenadas de 0 a 1 (\cref{SEC:DES-ANNOTATION}). Enviar esa tinta al reconocimiento habría
producido una imagen en blanco y, por tanto, un resultado vacío sin ningún error (gap-9).
Por eso cada reconocimiento usa un panel de escritura propio que genera la anotación en el
formato que el renderizador entiende: una anotación con \texttt{problem\_id} e
\texttt{is\_grade\_annotation} para las notas, y una anotación temporal, que se borra tras
leerla, para el texto. Cuando el servidor no aplica la nota por falta de confianza, la
interfaz ofrece usar el valor leído, que pasa entonces por el guardado manual con su
\texttt{old\_score}.

El cambio tuvo un efecto lateral que la revisión detectó antes de que llegara a producirse:
la carga de la tinta tomaba la primera anotación de trazo de cada página, de modo que una
anotación de nota, que se crea asociada a la página 1, habría sustituido a la tinta
de esa página tanto en el corrector como en la vista del estudiante. La carga ignora ahora
las anotaciones de los paneles.

\subsection{Administración, cuenta y notificaciones}

Esta subsección resume la implementación de la línea de gestión, desarrollada por Codex en
su propio \dfn{worktree} a partir de un borrador suyo revisado e integrado en esta memoria.
Comprende cuatro módulos: gestión de la organización (usuarios, cursos, configuración y
registro de auditoría), gestión de cada asignatura (miembros, grupos y permisos), perfil
personal y notificaciones.

\paragraph{Fidelidad al contrato.} Varias decisiones se tomaron porque el contrato no
autoriza más. El listado de usuarios solo declara parámetros de paginación, así que el
filtro de texto se aplica a la página recibida y la interfaz lo presenta como tal; el
registro de auditoría, en cambio, admite filtros remotos por actor, evento, entidad,
dirección IP e intervalo, que se envían al servidor. Aunque los cursos devuelven el campo
\texttt{is\_active}, el esquema de actualización solo admite etiqueta y fechas, por lo que no
se ofrece una acción de activación: la presencia de un campo en la respuesta no autoriza a
escribirlo. El perfil limita la edición al correo y a la preferencia de notificaciones,
únicos campos que admite su operación. Como no existe un listado de cursos, el catálogo se
obtiene del buscador jerárquico (\texttt{/search/}); esa carencia y la imposibilidad de
descubrir la membresía propia se registraron como huecos (gap-3 y gap-4).

\paragraph{Importaciones y exportaciones.} Las importaciones de usuarios y asignaturas
envían el fichero original en un \texttt{FormData} bajo el campo \texttt{file}, sin
convertirlo a un formato propio, y presentan los contadores de la respuesta, que puede
combinar filas correctas y errores. Los diagnósticos de error se tratan como objetos opacos,
porque el contrato no fija su estructura, y pueden descargarse en \acs{json}. Las
exportaciones distinguen la descarga binaria de usuarios de la respuesta estructurada de
asignaturas.

\paragraph{Alta, edición y baja de asignaturas.} El directorio de asignaturas
(\texttt{/admin/subjects}) permite al gestor crear, editar y eliminar asignaturas con
\texttt{POST /subjects/}, \texttt{PATCH /subjects/\{id\}/} y \texttt{DELETE /subjects/\{id\}/}.
El cuerpo de alta lleva nombre, código, identificador del coordinador y, solo si se rellena,
semestre; no lleva curso, porque la operación crea la asignatura en el curso activo de la
organización. Por eso el formulario advierte de ese destino aunque el directorio muestre un
curso histórico, cuyas filas tienen la edición y la baja deshabilitadas, y tras crear vuelve
al curso actual. La baja exige confirmación y conserva el destinatario si falla; tras una
mutación correcta el foco pasa al botón de creación, porque la fila que abrió el diálogo
puede haber desaparecido, y al borrar la última fila de una página posterior a la primera se
retrocede una página sin perder la búsqueda. Los manejadores simulados devuelven
\texttt{SubjectResponse} con los estados declarados (201, 200 y 204) y comprueban la unicidad
del código en el curso, la organización del coordinador y que la asignatura no tenga
exámenes antes de borrarla.

\paragraph{Permisos por asignatura.} La gestión de una asignatura se divide en pestañas
independientes para que cada una exija solo el permiso de su operación: quien puede
gestionar grupos puede usar esa pestaña aunque no pueda listar miembros. El editor de
permisos muestra, para cada permiso, el valor por defecto del rol y el efectivo, y ofrece
tres opciones (heredar, permitir, denegar); al guardar envía solo las excepciones. Durante el
desarrollo se detectó que, tras guardar una excepción y volver a \emph{heredar}, el
formulario recuperaba el valor anterior; se corrigió separando el estado de las excepciones
del de los permisos efectivos y se añadió una prueba de regresión.

\paragraph{Notificaciones.} El contador de la barra superior consulta los avisos sin leer
cada sesenta segundos mientras la pestaña está visible, se actualiza al volver a ella y
retira temporizador y escuchadores al desmontarse. Cada respuesta se vincula a la sesión que
la pidió y a un número de secuencia, de modo que una respuesta tardía no sustituye a una más
reciente ni muestra el contador de otra cuenta. Un fallo de consulta se muestra como contador
no disponible, no como cero.

\paragraph{Operaciones simultáneas y recuperación de errores.} El estado de React no basta
para impedir un doble envío: dos eventos atendidos antes del siguiente renderizado verían
ambos el formulario como libre. Por eso los formularios, las importaciones y las
confirmaciones usan, además del estado que deshabilita los botones, una referencia mutable
que se activa de forma síncrona antes de la petición y se libera en \texttt{finally}. La
misma referencia impide cerrar el diálogo, también con Escape, mientras la operación sigue
pendiente. Tras un fallo se conservan los valores introducidos y el fichero seleccionado, el
error se asocia al campo con \texttt{aria-describedby} y \texttt{aria-invalid} y el foco pasa
al primer campo inválido o al mensaje general; cuando una consulta fallida se recupera, el
foco vuelve a un destino concreto de la página (el buscador, el filtro de evento o la lista
de resultados), sin perder filtros ni página. Estas comprobaciones se hicieron sobre un
\acs{dom} simulado y no sustituyen una evaluación con tecnologías de apoyo
(\cref{SEC:AN-RNF}).
% Decisión editorial (19/10/2026): se mantiene este nivel de resumen; el borrador
% completo de la línea de gestión (unas 330 líneas) queda en mission/memoria-codex.
